\clearpage
\fancyhead[L]{\small BEOBACHTUNGEN.}
\SourceLeaf{468}
\begin{center}
\rule{0.55\linewidth}{0.4pt}\\[2pt]
Astronomische Nachrichten. Band X. Nr. 225. Seite 143. 3. Juli 1832.\\[2pt]
\rule{0.55\linewidth}{0.4pt}
\end{center}

\begin{center}[\textsc{Gauss} an \textsc{Schumacher.}]\quad Göttingen, 12. Mai 1832.\end{center}

{\rightskip=0pt plus 1em Beim letzten Merkurdurchgang habe ich die erste Berührung beim Eintritt wegen einer Wolke verpasst und die zweite durch einen Zufall. Den Austritt habe ich dagegen beobachtet\par}

\begin{center}
\begin{tabular}{ll}
innere Berührung beim Austritt \mbox{.\,.\,.\,.\,.\,.\,.} & $4^{\rm u}25^{\rm m}32^{\rm s}1$ \\
äußere Berührung \mbox{.\,.\,.\,.\,.\,.\,.\,.\,.\,.\,.\,.\,.\,.} & $4\;28\;22.5$ \\
\end{tabular}
\end{center}

Außerdem habe ich seine Kulmination am Meridiankreis beobachtet

\medskip
\noindent 5. Mai. Rektasz.\quad $2^{\rm u}49^{\rm m}50^{\rm s}36$\par
\noindent $\delta+16^\circ28'43^{\rm s}6$ der Nullpunkt aus dem Nadir bestimmt.\par
\noindent\hspace*{1em}---\quad ---\quad $44.4$ der Nullpunkt aus $\alpha$ Can. min. mit der Dekl. von $\alpha$ aus \textsc{Enckes} \mbox{I. B.}

